\section{Parallelism, perpendicularity, and angles}

The normal vector is the most efficient tool for comparing planes. A plane may be visually hard to compare with another plane, but its normal vector turns many geometric questions into vector questions.

Let
\[
\pi_1:A_1x+B_1y+C_1z+D_1=0,
\]
and
\[
\pi_2:A_2x+B_2y+C_2z+D_2=0.
\]
Their normal vectors are
\[
n_1=(A_1,B_1,C_1),
\qquad
n_2=(A_2,B_2,C_2).
\]

\subsection*{Parallel planes}

Two planes are parallel when their normal vectors are parallel. This means that one normal vector is a scalar multiple of the other.

\begin{definitionbox}
The planes \(\pi_1\) and \(\pi_2\) are parallel if there exists a non-zero scalar \(\lambda\) such that
\[
n_2=\lambda n_1.
\]
\end{definitionbox}

\begin{examplebox}
Consider
\[
\pi_1:2x-y+3z-1=0
\]
and
\[
\pi_2:4x-2y+6z+5=0.
\]
Their normal vectors are
\[
n_1=(2,-1,3),
\qquad
n_2=(4,-2,6).
\]
Since
\[
n_2=2n_1,
\]
the planes are parallel.
\end{examplebox}

Parallel planes may be distinct, or they may be the same plane. To check whether they are the same, test whether a point from one plane satisfies the equation of the other.

\subsection*{Perpendicular planes}

Two planes are perpendicular when their normal vectors are perpendicular.

\begin{definitionbox}
The planes \(\pi_1\) and \(\pi_2\) are perpendicular if
\[
n_1\cdot n_2=0.
\]
\end{definitionbox}

\begin{examplebox}
Consider
\[
\pi_1:x+2y-z=0
\]
and
\[
\pi_2:2x-y+0z+3=0.
\]
The normal vectors are
\[
n_1=(1,2,-1),
\qquad
n_2=(2,-1,0).
\]
Their dot product is
\[
n_1\cdot n_2=1\cdot2+2\cdot(-1)+(-1)\cdot0=0.
\]
Therefore the planes are perpendicular.
\end{examplebox}

The test is simple because the normal vectors carry the orientation of the planes.

\subsection*{Angle between planes}

The angle between two planes is measured by the angle between their normal vectors. Usually we take the acute angle, so we use an absolute value in the formula.

\begin{theorembox}
Let \(\theta\) be the acute angle between two planes with normal vectors \(n_1\) and \(n_2\). Then
\[
\cos\theta=\frac{|n_1\cdot n_2|}{\lVert n_1\rVert\,\lVert n_2\rVert}.
\]
\end{theorembox}

\begin{examplebox}
Determine the angle between
\[
\pi_1:x+2y-2z+1=0
\]
and
\[
\pi_2:2x-y+z-3=0.
\]
The normal vectors are
\[
n_1=(1,2,-2),
\qquad
n_2=(2,-1,1).
\]
Their dot product is
\[
n_1\cdot n_2=1\cdot2+2\cdot(-1)+(-2)\cdot1=2-2-2=-2.
\]
Their lengths are
\[
\lVert n_1\rVert=\sqrt{1^2+2^2+(-2)^2}=3,
\]
and
\[
\lVert n_2\rVert=\sqrt{2^2+(-1)^2+1^2}=\sqrt6.
\]
Therefore
\[
\cos\theta=\frac{|-2|}{3\sqrt6}=\frac{2}{3\sqrt6}.
\]
So
\[
\theta=\arccos\frac{2}{3\sqrt6}.
\]
\end{examplebox}

The calculation is a direct use of the dot product, but the interpretation is geometric: we compare the orientations of the planes through their normal directions.

\begin{summarybox}
When comparing two planes, ask:
\begin{itemize}
\item What are their normal vectors?
\item Are the normal vectors parallel?
\item Are the normal vectors perpendicular?
\item If an angle is needed, should it be the acute angle?
\item Do proportional normal vectors mean distinct parallel planes or the same plane?
\end{itemize}
\end{summarybox}

Plane comparison becomes much simpler when the normal vectors are read first.